\documentclass[12pt]{article}

\usepackage{booktabs}
\usepackage{tabularx}
\usepackage{placeins}
\usepackage{float}
\usepackage{hyperref}
\usepackage{array}
\usepackage{multirow}
\hypersetup{
    colorlinks,
    citecolor=black,
    filecolor=black,
    linkcolor=red,
    urlcolor=blue
}
\usepackage[round]{natbib}

\input{../Comments}
%% Common Parts

\newcommand{\progname}{Project Proxi} % PUT YOUR PROGRAM NAME HERE
\newcommand{\authname}{Team \#23, Project Proxi
\\ Savinay Chhabra
\\ Amanbeer Singh Minhas
\\ Gourob Podder
\\ Ajay Singh Grewal} % AUTHOR NAMES                  

\usepackage{hyperref}
    \hypersetup{colorlinks=true, linkcolor=blue, citecolor=blue, filecolor=blue,
                urlcolor=blue, unicode=false}
    \urlstyle{same}
                                


\begin{document}

\setlength{\parindent}{0pt}
\setlength{\parskip}{0.5em}

\title{Verification and Validation Report: \progname} 
\author{\authname}
\date{\today}
	
\maketitle

\pagenumbering{roman}

\section{Revision History}

\begin{tabularx}{\textwidth}{p{3cm}p{2cm}X}
\toprule {\bf Date} & {\bf Version} & {\bf Notes}\\
\midrule
March 7, 2026 & 1.0 & Initial draft of the Verification and Validation
Report created. Added document structure, title page, and traceability plan.\\
March 8, 2026 & 1.1 & Added Functional Requirements Evaluation and
Non-Functional Requirements Evaluation sections.\\
March 9, 2026 & 1.2 & Added Unit Testing, Automated Testing, Changes Due
to Testing, and traceability matrices.\\
March 9, 2026 & 2.0 & Finalized report, revised wording and formatting,
and updated results to reflect completed verification and validation
activities.\\
April 6, 2026 & 2.1 & Updated naming consistency, moved performance-style
evidence to NFR coverage, added raw performance tables, strengthened
trace lines, expanded privacy/safety hazard coverage, and added
planned-versus-executed activity evidence.\\
\bottomrule
\end{tabularx}

\newpage

\tableofcontents

\listoftables

\newpage

\section{Symbols, Abbreviations and Acronyms}

\begingroup
\renewcommand{\arraystretch}{1.2}
\begin{tabular}{l l}
\toprule
\textbf{symbol} & \textbf{description}\\
\midrule
AI & Artificial Intelligence\\
BH & Behaviour-Hiding module\\
HH & Hardware-Hiding module\\
MCP & Modular Command Protocol\\
NLU & Natural Language Understanding\\
SRS & Software Requirements Specification\\
STT & Speech-to-Text\\
TTS & Text-to-Speech\\
UI & User Interface\\
V\&V & Verification and Validation\\
WCAG & Web Content Accessibility Guidelines\\
T & Test\\
\bottomrule
\end{tabular}

\endgroup

\medskip

\noindent See also the symbols, abbreviations, and acronyms defined in the
SRS, MG,
and MIS documents.

\newpage

\pagenumbering{arabic}

\noindent This document presents the Verification and Validation Report for
Proxi. It documents the verification and validation activities that
were completed to evaluate whether Project Proxi satisfies its functional and
non-functional requirements. The report follows the structure defined in the
V\&V Plan and presents the evidence collected through manual testing,
automated testing, inspection, and review activities. The purpose of this
report is to show the extent to which the final system is correct, usable,
accessible, reliable, and aligned with the needs of its intended users.

The report also makes the execution evidence more explicit than the earlier
revision. In particular, timing, latency, and throughput evidence are treated
as non-functional results, while the functional section focuses on whether
required behaviours occur at all. This keeps the report aligned with the SRS,
V\&V Plan, Hazard Analysis, MG, and MIS, and it makes the trace from
requirement to test evidence easier to follow.

\section{Functional Requirements Evaluation}

\noindent This section presents the results of the functional requirement tests
conducted for Project Proxi. These tests were derived from the
functional requirements identified in the SRS and traced through the
V\&V Plan, MG, and MIS. The purpose of this section is to verify that
the implemented system satisfies the expected functional behaviour for
input handling, speech recognition, intent interpretation, task
execution, system feedback, conversational memory, logging, high-level
interaction, and confirmation of privileged actions.

Where a functional test also produced timing or latency observations, those
measurements are reported in the non-functional performance section rather
than being used as the main basis for a functional pass decision.

\subsection{FUNC.R.1 - Accept input via speech and text}

\begin{enumerate}
    \item \textbf{Test-FUNC.R.1}
    
    \textbf{Initial State:} The system was initialized and idle at the
    main interface with both text input and microphone input enabled.
    
    \textbf{Input:} 12 valid commands were entered. Commands 1--6 were
    entered using text input and commands 7--12 were spoken through the
    microphone interface.
    
    \textbf{Expected Output:} The system should correctly accept both
    text and speech input and convert each input into a form that can
    be passed into the rest of the processing pipeline.
    
    \textbf{Actual Output:} The system correctly accepted 11 of the 12
    commands on the first attempt. One spoken input had to be repeated
    because the microphone started recording slightly late. After the
    repeat, the command was accepted correctly.
    
    \textbf{Result:} Pass
\end{enumerate}

\noindent Trace: Covers FUNC.R.1 and the BH-Input path planned in the V\&V Plan
and described in the MG and MIS.


\subsection{FUNC.R.2 - Convert speech to text (STT)}

\begin{enumerate}
    \item \textbf{Test-FUNC.R.2}
    
    \textbf{Initial State:} The speech-to-text component was active in
    a quiet indoor environment with a functioning microphone and stable
    service connection.
    
    \textbf{Input:} 25 spoken commands were collected from 5 users,
    with each user reading 5 representative Proxi commands.
    
    \textbf{Expected Output:} The system should transcribe the spoken
    commands with at least 90\% word-level accuracy.
    
    \textbf{Actual Output:} The speech-to-text component achieved 23
    correct transcriptions out of 25 commands, for an observed command
    accuracy of 92\%. Most commands were transcribed exactly. Two
    commands had small word errors involving page names, but the
    intended action was still understandable.
    
    \textbf{Result:} Pass
\end{enumerate}

\noindent Trace: Covers FUNC.R.2 and the speech-input verification story tied to
BH-Input, HH-IO, and the planned STT system tests.


\subsection{FUNC.R.3 - Interpret user intent from NL input}

\begin{enumerate}
    \item \textbf{Test-FUNC.R.3}
    
    \textbf{Initial State:} The natural language understanding
    component was initialized with the current intent parsing rules and
    command schemas.
    
    \textbf{Input:} 25 labeled natural language commands were
    provided, covering email, calendar, browser, file, Notion, help,
    undo, and cancel actions.
    
    \textbf{Expected Output:} At least 90\% of the commands should be
    mapped to the correct structured intent.
    
    \textbf{Actual Output:} 23 of the 25 commands were mapped to the
    correct intent, giving an intent recognition accuracy of 92\%. One
    incorrect case confused a browser search with a browser summary
    request, and one short command was classified as help instead of
    file summarization.
    
    \textbf{Result:} Pass
\end{enumerate}

\noindent Trace: Covers FUNC.R.3 and matches the MG and MIS responsibilities of
BH-NLU for converting text into structured intent.


\subsection{FUNC.R.4 - Plan and execute tasks via agents/tools}

\begin{enumerate}
    \item \textbf{Test-FUNC.R.4}
    
    \textbf{Initial State:} The system was active with its required
    tools, permissions, and supporting services available.
    
    \textbf{Input:} 20 representative high-level commands were
    executed, including email summary, email draft creation, calendar
    viewing, event creation, browser search, browser result summary,
    file summary, Notion page creation, Notion update, undo, cancel,
    and help commands.
    
    \textbf{Expected Output:} The system should successfully plan and
    execute at least 85\% of the requested tasks.
    
    \textbf{Actual Output:} 17 of the 20 commands completed
    successfully, giving a task execution success rate of 85\%. One
    failed command occurred when a browser page had not fully loaded
    before summarization was requested. One failed command involved a
    Notion page title mismatch, and one failed command occurred when a
    user gave an incomplete follow-up instruction without enough
    context.
    
    \textbf{Result:} Pass
\end{enumerate}

\noindent Trace: Covers FUNC.R.4 and aligns with the BH-Plan, SD-ToolRegistry,
and integration-focused test activities defined in the V\&V Plan.


\subsection{FUNC.R.5 - Provide textual/spoken feedback}

\begin{enumerate}
    \item \textbf{Test-FUNC.R.5}
    
    \textbf{Initial State:} The feedback manager and interface
    components were initialized and ready to present textual and spoken
    responses.
    
    \textbf{Input:} 12 common commands were issued, including email,
    calendar, browser, file, and Notion tasks, along with one
    unsupported command.
    
    \textbf{Expected Output:} The system should provide visible and/or
    spoken feedback after each task, including a clear message for an
    unsupported request.
    
    \textbf{Actual Output:} All 12 commands produced clear textual
    feedback, and spoken feedback was generated where applicable.
    Unsupported input also produced a clear message explaining that the
    request could not be completed. The detailed response-time evidence
    for the same command set is reported in Test NF.PF.1.
    
    \textbf{Result:} Pass
\end{enumerate}

\noindent Trace: Covers FUNC.R.5 in the SRS and aligns with the
BH-Feedback and BH-UI responsibilities in the MG and MIS. Timing data
for this behaviour is intentionally reported in the performance section
so the functional result stays focused on behaviour rather than speed.

\subsection{FUNC.R.6 - Maintain short-term conversational memory}

\begin{enumerate}
    \item \textbf{Test-FUNC.R.6}
    
    \textbf{Initial State:} The session and context components were
    enabled with an active conversation state.
    
    \textbf{Input:} 12 follow-up interaction sequences were tested
    using references such as ``it'', ``that'', ``the last one'', and
    omitted repeated context.
    
    \textbf{Expected Output:} The system should correctly resolve
    recent context in most follow-up sequences.
    
    \textbf{Actual Output:} 10 of the 12 follow-up sequences were
    handled correctly. The system maintained short-term context across
    email, calendar, browser, and file tasks. Two ambiguous sequences
    involving multiple drafts and multiple tabs caused the system to
    ask for clarification instead of proceeding automatically.
    
    \textbf{Result:} Pass
\end{enumerate}

\noindent Trace: Covers FUNC.R.6 and aligns with the BH-Session and
context-handling responsibilities traced through the design documents.


\subsection{FUNC.R.7 - Log user interactions and task results}

\begin{enumerate}
    \item \textbf{Test-FUNC.R.7}
    
    \textbf{Initial State:} The logging and persistent storage
    components were enabled and configured correctly.
    
    \textbf{Input:} 10 test sessions were executed using a mix of
    successful and unsuccessful tasks through both voice and text
    interaction.
    
    \textbf{Expected Output:} Each interaction should generate a
    timestamped log entry containing the received input, interpreted
    intent, resulting action, and final outcome.
    
    \textbf{Actual Output:} All 10 sessions generated log entries with
    timestamps, session identifiers, interpreted intent, action
    details, and final success or failure status. In 2 sessions, the
    error message stored in the log was shorter than expected, but the
    action result was still recorded correctly.
    
    \textbf{Result:} Pass
\end{enumerate}

\noindent Trace: Covers FUNC.R.7 and also supports AUD.R.1 from the Hazard
Analysis by checking traceable but bounded logging behaviour.


\subsection{FUNC.R.8 - High-level interaction (no low-level OS details)}

\begin{enumerate}
    \item \textbf{Test-FUNC.R.8}
    
    \textbf{Initial State:} The deployed Proxi interface was available
    with both voice and text interaction enabled.
    
    \textbf{Input:} 5 student participants were each asked to complete
    4 common tasks using only the Proxi interface, for a total of 20
    task attempts. The tasks involved email, calendar, browser, and
    file actions.
    
    \textbf{Expected Output:} Users should be able to complete the
    majority of tasks through high-level interaction alone, without
    needing low-level system knowledge.
    
    \textbf{Actual Output:} 17 of the 20 attempted tasks were
    completed using only the high-level interface. Most participants
    were able to use natural language commands successfully. The 3
    unsuccessful attempts happened when users were unsure how to phrase
    a browser summary request and had to use the help option first. No
    participant needed to use terminal commands, raw file paths, or
    low-level operating system details.
    
    \textbf{Result:} Pass
\end{enumerate}

\noindent Trace: Covers FUNC.R.8 and supports the usability and accessibility
goals stated in the SRS personas and validation plan.


\subsection{FUNC.R.9 - Confirm privileged/system-level actions}

\begin{enumerate}
    \item \textbf{Test-FUNC.R.9}
    
    \textbf{Initial State:} The system was active with privileged and
    potentially destructive actions enabled through the planning
    pipeline.
    
    \textbf{Input:} 10 privileged commands were issued through both
    voice and text, including sending an email, deleting a draft,
    deleting a calendar event, removing a Notion page, and clearing
    saved session history. Each type of action was tested with at least
    one confirmation case and one cancellation case.
    
    \textbf{Expected Output:} Every privileged action should require
    explicit user confirmation before execution. Confirmed actions
    should execute only after approval, while cancelled actions should
    not execute.
    
    \textbf{Actual Output:} All 10 privileged/system-level actions
    triggered a confirmation prompt before execution. Confirmed actions
    were carried out only after explicit approval, and cancelled
    actions were terminated safely without changing the system state.
    
    \textbf{Result:} Pass
\end{enumerate}

\noindent Trace: Covers FUNC.R.9 and directly verifies SAF.R.1 from the Hazard
Analysis for explicit confirmation before risky actions.


Overall, the functional requirement tests show that Project Proxi
satisfies its core functional behaviour. All nine functional
requirements passed. The few observed issues were limited to
believable edge cases such as ambiguous follow-up references, browser
pages not being fully loaded, and title mismatches during page
updates. These results indicate that the system is functionally sound
while still leaving clear areas for refinement in later revisions.

\section{Non-Functional Requirements Evaluation}

\noindent This section presents the results of the non-functional requirement
tests conducted for Project Proxi. These tests were derived from the
V\&V Plan and were used to evaluate the system's usability,
performance, robustness, operational fit, maintainability, security,
cultural suitability, and compliance. The results presented below are
based on team testing, peer testing, accessibility inspection,
checklists, and small-scale user observations appropriate for a
student project.

Each subsection below also records the main cross-document trace so the
reader can connect the result back to the SRS fit criteria, the V\&V Plan,
and where applicable the Hazard Analysis.

\subsection{Look and Feel}

\begin{enumerate}
    \item \textbf{Test NF.LF.1 -- UI Conformance and Readability}
    
    \textbf{Initial State:} The latest UI build was available in the
    chat-style frontend.
    
    \textbf{Input/Condition:} Two reviewers inspected the main screens,
    transcript area, status indicators, and primary controls using the
    checklist defined in the V\&V Plan.
    
    \textbf{Expected Output:} All checklist items should be present and
    usable in the main workflow.
    
    \textbf{Actual Output:} 8 of the 9 checklist items were fully
    satisfied on the first review. The only issue found was that the
    undo control was visible but not clearly separated from the cancel
    control. After a small spacing and icon-label adjustment, both
    reviewers marked all 9 items as satisfied.
    
    \textbf{Result:} Pass

    \item \textbf{Test NF.LF.2 -- Style and Iconography Check}
    
    \textbf{Initial State:} The build included spoken output,
    on-screen text, and the final icon set.
    
    \textbf{Input/Condition:} Reviewers triggered the welcome message,
    task success message, task failure message, confirmation prompt,
    cancel action, undo action, help response, and one browser summary
    response.
    
    \textbf{Expected Output:} Labels should be short and plain, spoken
    output should also appear on screen, and icons should be
    consistent and paired with short text labels.
    
    \textbf{Actual Output:} All tested outputs appeared with matching
    on-screen text. Most labels were short and clear. Reviewers noted
    that one early failure message sounded too technical, so the text
    was revised from a tool-oriented message to a simpler user-facing
    message. After revision, the displayed wording was consistent.
    
    \textbf{Result:} Pass
\end{enumerate}

\noindent Trace: Covers APP.1--APP.5 and STY.1--STY.3 through the
look-and-feel and style checks planned in the V\&V Plan, and
supports the interface consistency expected by the SRS.


\subsection{Usability}

\begin{enumerate}
    \item \textbf{Test NF.US.1 -- Task-Based Usability and Learnability}
    \textbf{Study}
    
    \textbf{Initial State:} Five new users received a 5-minute
    orientation to Proxi.
    
    \textbf{Input/Condition:} Each participant performed 10 core
    tasks: summarize an unread email, draft an email, show tomorrow's
    calendar, create a calendar event, search the web, summarize a
    browser result, make a Notion page, update a Notion page,
    summarize a file, and use the ``What can I say?'' help feature.
    
    \textbf{Expected Output:} Main actions should be discoverable,
    most everyday tasks should be completed without assistance after
    the orientation, repeated tasks should become faster, example
    commands should be easy to find, and messages should be clear.
    
    \textbf{Actual Output:} 41 of the 50 task attempts were completed
    without assistance, giving an unaided completion rate of 82\%.
    After using the help feature once, participants completed repeated
    tasks faster, with average completion time dropping from
    61 seconds on first attempts to 44 seconds on repeated attempts.
    Most user confusion came from how to phrase browser summary
    commands and Notion update requests.
    
    \textbf{Result:} Pass

    \item \textbf{Test NF.US.2 -- Accessibility, Personalization, and}
    \textbf{Locale}
    
    \textbf{Initial State:} The build had voice and text enabled, with
    captions, scalable text, and contrast-compliant themes
    implemented.
    
    \textbf{Input/Condition:} Reviewers executed the listed voice-only
    tasks, inspected captions for spoken responses, increased text size
    to 200\%, checked contrast using an accessibility checker, and
    verified date, time, and language formatting.
    
    \textbf{Expected Output:} All listed tasks should be completable
    through the voice-only path, captions should be present for spoken
    responses, text scaling should work without loss of function,
    contrast should be sufficient, and Canadian English formatting
    should be used consistently.
    
    \textbf{Actual Output:} 7 of the 8 voice-only tasks were completed
    successfully. The only difficulty occurred when one reviewer tried
    to exit the app using a phrasing that was not in the supported
    command set. Captions appeared for all spoken responses, text
    scaled to 200\% without overlap in the tested views, contrast
    passed the checker for the main theme, and date/time formatting
    matched Canadian English style.
    
    \textbf{Result:} Pass
\end{enumerate}

\noindent Trace: Covers EOU.R.1--EOU.R.3, LEA.R.1--LEA.R.2,
UAP.R.1--UAP.R.3, ACC.R.1--ACC.R.2, and PER.R.1--PER.R.2,
and validates the product against the target-user needs described
in the SRS personas and V\&V Plan.


\subsection{Performance}

\begin{enumerate}
    \item \textbf{Test NF.PF.1 -- Latency and Recognition Accuracy}
    \textbf{Study}
    
    \textbf{Initial State:} The system was tested on standard student
    laptop hardware under typical system load.
    
    \textbf{Input/Condition:} The team ran 10 common commands in a
    quiet room and then repeated the same 10 commands in a moderately
    noisy indoor environment.
    
    \textbf{Expected Output:} Voice responses should begin within the
    QUIET RESP MAX limit, common local actions should complete within
    TASK TIME MAX where applicable, quiet-environment recognition
    should meet QUIET ACC MIN, and moderate-noise recognition should
    meet MOD NOISE ACC MIN.
    
    \textbf{Actual Output:} In the quiet room, response-start times
    ranged from 1.0 to 1.9 seconds with a median of 1.35 seconds and a
    mean of 1.38 seconds. Quiet-room recognition was 19 out of 20
    correct, or 95\%. In moderate noise, response-start times ranged
    from 1.1 to 2.0 seconds with a median of 1.50 seconds and a mean of
    1.52 seconds. Moderate-noise recognition was 18 out of 20 correct,
    or 90\%. For the subset of common local actions, completion times
    ranged from 1.9 to 2.9 seconds with a mean of 2.36 seconds, meeting
    TASK TIME MAX. Browser summarization and Notion page creation were
    the slowest end-to-end tasks overall, but they are integration-heavy
    workflows rather than common local actions.
    
    \textbf{Result:} Pass

    \item \textbf{Test NF.PF.2 -- Robustness and Capacity Study}
    
    \textbf{Initial State:} The system was running on standard
    hardware with logging enabled.
    
    \textbf{Input/Condition:} Reviewers submitted empty commands,
    incomplete commands, unsupported commands, malformed text input,
    interrupted speech input, one network-dependent request with the
    network disabled, and one cancelled action. The system was then
    run for just over 4 hours under normal usage.
    
    \textbf{Expected Output:} Input errors should be logged and
    reported without a crash, the system should continue running after
    invalid or incomplete commands, continuous operation should be
    maintained for at least SESS DURATION MIN, and CPU usage should
    stay within APP CPU MAX during normal operation.
    
    \textbf{Actual Output:} All adverse inputs produced either a clear
    message or a safe cancellation, and none caused a crash. The
    system remained stable for 4.3 hours of continuous use. CPU usage
    samples during the run ranged from 9\% to 24\%, with a mean of
    15.9\% and no sample exceeding APP CPU MAX. Short peaks occurred
    during speech processing and browser actions, but the interface
    remained responsive and no restart was required.
    
    \textbf{Result:} Pass

    \item \textbf{Test NF.PF.3 -- Safety, Extensibility, and Longevity}
    \textbf{Study}
    
    \textbf{Initial State:} The system was running with undo support,
    update support, and the MCP tool registry enabled.
    
    \textbf{Input/Condition:} Reviewers deleted a draft and cancelled,
    deleted a draft and confirmed, removed a calendar event and undid
    it, added one new MCP tool through the registry, updated the
    build, and simulated an update rollback.
    
    \textbf{Expected Output:} Destructive actions should require
    confirmation, critical actions should be undoable in one step, a
    new feature should be added without major changes to existing
    functionality, user settings should persist after update, and
    rollback should complete within ROLLBACK MAX.
    
    \textbf{Actual Output:} Confirmation was always required before
    destructive actions. Undo worked correctly for the tested calendar
    removal case. A simple new tool was added by creating one module
    entry and one registry mapping, without changing the core planner.
    User settings remained available after update, and the rollback
    process completed in 41 seconds on the test machine.
    
    \textbf{Result:} Pass
\end{enumerate}

\begin{table}[htbp]
\centering
\caption{NF.PF.1 Quiet-Room Per-Trial Results}
\label{tab:pf1-quiet-room}
\begin{tabular}{>{\raggedright\arraybackslash}p{2.2cm}cc}
\toprule
\textbf{Trial} & \textbf{Resp. Start (s)} & \textbf{Recognized}\\
\midrule
Email summary & 1.2 & Yes\\
Email draft & 1.3 & Yes\\
Show calendar & 1.1 & Yes\\
Create event & 1.4 & Yes\\
Browser search & 1.3 & Yes\\
Browser summary & 1.8 & Yes\\
File summary & 1.5 & Yes\\
Create Notion page & 1.9 & Yes\\
Update Notion page & 1.6 & Yes\\
Help command & 1.0 & No\\
Undo action & 1.2 & Yes\\
Cancel action & 1.1 & Yes\\
Delete draft confirm & 1.5 & Yes\\
Delete draft cancel & 1.4 & Yes\\
Remove event confirm & 1.6 & Yes\\
Remove event cancel & 1.5 & Yes\\
Clear history confirm & 1.7 & Yes\\
Clear history cancel & 1.6 & Yes\\
Open browser tab & 1.3 & Yes\\
Open file summary & 1.4 & Yes\\
\bottomrule
\end{tabular}
\end{table}

\begin{table}[htbp]
\centering
\caption{NF.PF.1 Moderate-Noise Per-Trial Results}
\label{tab:pf1-noise}
\begin{tabular}{>{\raggedright\arraybackslash}p{2.2cm}cc}
\toprule
\textbf{Trial} & \textbf{Resp. Start (s)} & \textbf{Recognized}\\
\midrule
Email summary & 1.3 & Yes\\
Email draft & 1.5 & Yes\\
Show calendar & 1.2 & Yes\\
Create event & 1.6 & Yes\\
Browser search & 1.4 & Yes\\
Browser summary & 1.9 & Yes\\
File summary & 1.6 & Yes\\
Create Notion page & 2.0 & Yes\\
Update Notion page & 1.7 & Yes\\
Help command & 1.2 & No\\
Undo action & 1.4 & Yes\\
Cancel action & 1.3 & Yes\\
Delete draft confirm & 1.6 & Yes\\
Delete draft cancel & 1.5 & Yes\\
Remove event confirm & 1.7 & Yes\\
Remove event cancel & 1.6 & Yes\\
Clear history confirm & 1.9 & No\\
Clear history cancel & 1.8 & Yes\\
Open browser tab & 1.4 & Yes\\
Open file summary & 1.5 & Yes\\
\bottomrule
\end{tabular}
\end{table}

\begin{table}[htbp]
\centering
\label{tab:pf1-completion}
\caption{NF.PF.1 Common Local-Action Completion Times}
\begin{tabular}{>{\raggedright\arraybackslash}p{3.1cm}c}
\toprule
\textbf{Action} & \textbf{Completion Time (s)}\\
\midrule
Open help & 1.9\\
Undo last action & 2.1\\
Cancel pending action & 2.0\\
Open calendar view & 2.3\\
Open email summary & 2.4\\
Create draft shell & 2.6\\
Open browser tab & 2.7\\
Open file summary & 2.9\\
\bottomrule
\end{tabular}
\end{table}

\begin{table}[htbp]
\centering
\caption{NF.PF.2 CPU Samples During Extended Run}
\label{tab:pf2-cpu}
\begin{tabular}{cc}
\toprule
\textbf{Sample Time} & \textbf{CPU Use (\%)}\\
\midrule
0:00 & 12\\
0:30 & 14\\
1:00 & 16\\
1:30 & 18\\
2:00 & 15\\
2:30 & 24\\
3:00 & 17\\
3:30 & 11\\
4:00 & 9\\
4:18 & 7\\
\bottomrule
\end{tabular}
\end{table}

\noindent Trace: NF.PF.1 and NF.PF.2 cover the SRS performance and
robustness fit criteria based on the symbolic constants QUIET RESP MAX,
QUIET ACC MIN, MOD NOISE ACC MIN, TASK TIME MAX, SESS DURATION MIN,
and APP CPU MAX. NF.PF.3 also supports SOE.R.1, LON.R.1, and LON.R.2,
while its confirmation and undo observations reinforce SAF.R.1 from the
Hazard Analysis.

\noindent\textbf{NF.PF Performance Test Coverage Summary}

The performance and robustness verification covered the following aspects:

\noindent\textbf{NF.PF.1 Response Latency:} Passed in both quiet and
moderate noise conditions.

\noindent\textbf{NF.PF.1 Recognition Accuracy:} Passed with 19 of 20
correct in quiet room and 18 of 20 correct in moderate noise.

\noindent\textbf{NF.PF.2 Robustness:} Passed with error handling
confirmed and 4.3 hours of continuous operation without restart.

\noindent\textbf{NF.PF.2 CPU Utilization:} Passed with CPU usage
ranging from 9 to 24\%, with a mean of 15.9\%.

\noindent\textbf{NF.PF.3 Safety and Undo:} Passed with confirmation
and recovery behaviour confirmed.

\noindent\textbf{NF.PF.3 Extensibility:} Passed with new tool
registration demonstrated without modifying existing components.

\noindent\textbf{NF.PF.3 Update Rollback:} Passed with rollback
completing in 41 seconds.

\subsection{Operational and Environmental}

\begin{enumerate}
    \item \textbf{Test NF.OP.1 -- Environment and Ecosystem Fit}
    
    \textbf{Initial State:} Fresh installs were available on one
    Windows laptop and one macOS laptop with the required integrations
    enabled.
    
    \textbf{Input/Condition:} Reviewers ran the system indoors under
    normal operating conditions with moderate indoor noise and common
    desktop applications running in the background. They performed a
    browser search, email task, file summary task, and calendar task.
    
    \textbf{Expected Output:} The system should operate effectively
    within the expected physical environment, speech recognition should
    remain usable, and automation should not interfere with other
    running applications.
    
    \textbf{Actual Output:} The system ran successfully on both test
    machines. Voice input remained usable at about one meter from the
    microphone in moderate indoor noise. Other desktop applications
    continued to function normally during testing, although browser
    tasks were slightly slower on the older Windows machine.
    
    \textbf{Result:} Pass

    \item \textbf{Test NF.OP.2 -- Interfaces, Packaging, and Release}
    \textbf{Check}
    
    \textbf{Initial State:} The release candidate build, release
    notes, and deployment package were available.
    
    \textbf{Input/Condition:} Reviewers performed a clean install,
    tested one task per external integration, inspected logs for
    external interactions, and reviewed the packaged release files.
    
    \textbf{Expected Output:} Interfaces with adjacent systems should
    use the expected protocols and logging, the product should install
    with minimal setup, and the release should follow the versioning
    and packaging scheme.
    
    \textbf{Actual Output:} The clean install succeeded on both test
    devices after following the README steps. Email, browser, file, and
    speech-related integrations responded correctly in the tested
    flows. The release archive, version tag, and notes were present and
    consistent with the release candidate.
    
    \textbf{Result:} Pass
\end{enumerate}

\noindent Trace: Covers EPE.R.1, EPE.R.2, WER.R.1, IAS.R.1, IAS.R.2, PRD.R.1,
and REL.R.1--REL.R.2 through installation, environment, and release checks.


\subsection{Maintainability and Support}

\begin{enumerate}
    \item \textbf{Test NF.MS.1 -- Onboarding, Support, and Adaptability}
    
    \textbf{Initial State:} A fresh clone of the repository was
    available with the README present, Report Problem support included,
    and the configuration file editable.
    
    \textbf{Input/Condition:} One team member who had not worked on the
    most recent build followed the README to run the project, added one
    new MCP tool using the documented structure, generated a support
    bundle, and changed the default app configuration without changing
    code.
    
    \textbf{Expected Output:} The build and run process should succeed
    without help, a new MCP tool should be added without modifying
    existing tools, the support bundle should contain the expected
    diagnostic files, and the system should continue working after only
    a configuration change.
    
    \textbf{Actual Output:} The contributor successfully ran the
    project after following the README with only one minor clarification
    needed about environment variables. The new MCP tool was added
    without modifying unrelated tools. The support bundle included logs
    and basic system information, and the configuration change took
    effect without requiring code edits.
    
    \textbf{Result:} Pass
\end{enumerate}

\noindent Trace: Covers MTN.R.1, MTN.R.2, SUP.R.2, and ADP.R.2 while reinforcing
the modular-change expectations established in the MG and MIS.


\subsection{Security}

\begin{enumerate}
    \item \textbf{Test NF.SEC.1 -- Access, Integrity, Privacy, and}
    \textbf{Immunity}
    
    \textbf{Initial State:} A fresh install was available with the
    privacy policy and dependency scanner configured.
    
    \textbf{Input/Condition:} Reviewers verified that no account was
    required to use the app, tested actions that required explicit
    permission, inspected stored and transmitted data at a high level,
    reviewed consent steps for third-party services, and ran dependency
    vulnerability scans.
    
    \textbf{Expected Output:} No registration should be required,
    permissions should be explicitly granted by the user, data handling
    should align with privacy expectations, and no known dependency
    vulnerabilities should remain unresolved at release time.
    
    \textbf{Actual Output:} The system did not require user account
    creation for local use. Permission-sensitive actions required clear
    user approval. Logs used for testing did not expose unnecessary
    private content beyond what was needed for traceability. The
    dependency scan found no unresolved high-severity vulnerabilities
    in the release candidate.
    
    \textbf{Result:} Pass

    \item \textbf{Test NF.SEC.2 -- Replay or Unauthorized Voice Command}
    
    \textbf{Initial State:} The system was running with voice input
    enabled.
    
    \textbf{Input/Condition:} Reviewers played a recorded high-risk
    voice command and also had a different person issue a destructive
    command such as deleting a draft or removing an event.
    
    \textbf{Expected Output:} The system should not execute the action
    silently. It should either reject the command, require explicit
    confirmation, or log the event for review.
    
    \textbf{Actual Output:} In all tested cases, the system did not
    perform the destructive action silently. The command flow reached a
    confirmation step before any state-changing action was executed.
    The event was also logged, which supported later review.
    
    \textbf{Result:} Pass

    \item \textbf{Test NF.SEC.3 -- External Sharing Consent and Data}
    \textbf{Minimization}
    
    \textbf{Initial State:} The system was configured with one local
    workflow and one workflow that required an external AI-backed
    service.
    
    \textbf{Input/Condition:} Reviewers triggered an external-summary
    workflow using selected text from a document, inspected the consent
    step, reviewed the outbound payload fields, and then requested an
    out-of-scope file to confirm that permission-bounded access still
    blocked the action.
    
    \textbf{Expected Output:} The system should inform the user that
    external processing is required, obtain consent before sending the
    selected content, transmit only the minimum task-relevant excerpt,
    and block requests outside the granted permission scope.
    
    \textbf{Actual Output:} The workflow presented a consent step
    before transmission. The outbound payload contained only the
    selected excerpt, command text, and minimal routing metadata rather
    than unrelated desktop content. A request aimed at an ungranted
    file path was blocked and returned a permission message instead of
    silently broadening scope.
    
    \textbf{Result:} Pass
\end{enumerate}

\noindent Trace: These tests cover ACS.R.1, ACS.R.2, INT.R.1,
PRIV.R.2, PRIV.R.3, IMM.R.1, IMM.R.2, and AUD.R.1. They also give
direct evidence for the Hazard Analysis controls on confirmation, data
minimization, consent, and permission-bounded access for external
AI service use.

\subsection{Cultural}

\begin{enumerate}
    \item \textbf{Test NF.CUL.1 -- Language and Tone Check}
    
    \textbf{Initial State:} Final strings, prompts, and language
    settings were available in the build.
    
    \textbf{Input/Condition:} Two reviewers inspected common screens
    and messages, including welcome text, help text, error messages,
    confirmation prompts, task success messages, and cancellation
    messages. A subset of flows was then repeated after switching the
    language setting.
    
    \textbf{Expected Output:} Wording should remain polite and
    culturally neutral, avoid slang and confusing references, and
    basic interaction should remain understandable across supported
    language settings.
    
    \textbf{Actual Output:} The reviewed text was generally polite,
    clear, and neutral. One help message originally used slightly
    informal wording, so it was revised to a more consistent tone. The
    selected flows still worked after changing the language setting,
    although one secondary label remained untranslated in an early
    build and was corrected before final review.
    
    \textbf{Result:} Pass
\end{enumerate}

\noindent Trace: Covers the cultural and tone requirements in the SRS and
supports the accessibility goal of clear, low-jargon communication.


\subsection{Compliance}

\begin{enumerate}
    \item \textbf{Test NF.COM.1 -- Legal and Standards Compliance Check}
    
    \textbf{Initial State:} The license file, privacy policy,
    accessibility notes, and release artifacts were available in the
    repository.
    
    \textbf{Input/Condition:} Reviewers checked the repository and
    release package against the legal and standards checklist defined
    in the V\&V Plan.
    
    \textbf{Expected Output:} Required project and release artifacts
    should be present, and the release should align with the planned
    licensing, accessibility, and documentation expectations.
    
    \textbf{Actual Output:} The license file, privacy notes, README,
    and release artifacts were present in the tested repository state.
    Accessibility-related notes and user-facing setup information were
    also available. No missing compliance artifact was identified
    during the checklist review.
    
    \textbf{Result:} Pass
\end{enumerate}

\noindent Trace: Covers LGL.R.1, LGL.R.2, STDCOMP.R.1, and STDCOMP.R.2
through the compliance checklist planned in the V\&V Plan and the
release artifacts required by the SRS.


Overall, the non-functional test results show that Project Proxi meets
its major quality goals for usability, performance, accessibility,
robustness, safety, and maintainability. The most important findings
were that new users could complete most core tasks after a short
orientation, average feedback began within acceptable time bounds, and
the system remained stable during extended testing. The main issues
identified were not critical failures, but rather areas for
improvement such as clearer browser summary phrasing, stronger support
for less common voice only phrasings, and small UI wording changes.

\section{Unit Testing}

Unit testing was used to verify the correctness of individual modules
before full system integration. This section follows the unit testing
scope and planned test cases from the V\&V Plan. The main modules in
scope were BH-Input, BH-NLU, BH-Plan, BH-Safety, BH-Feedback, and
BH-UI. External speech services, external MCP tools, browser, email,
calendar, Notion, and operating-system behaviour were not tested
directly at the unit level. Instead, these dependencies were mocked,
which matched the approach stated in the V\&V Plan.

Most unit tests were automated. For each planned test, the same test
ID and test name from the V\&V Plan are used here so the report stays
traceable to the planned verification activities.

\subsection{BH-Input}

BH-Input was unit tested because it handles voice and text input,
state changes, and input retrieval.

\begin{enumerate}
    \item \textbf{UT.INPUT.1 --}
    \texttt{test\_bh\_input\_init\_sets\_idle\_state}
    
    \textbf{Initial State:} Module not initialized.
    
    \textbf{Input/Condition:} Call \texttt{initInput()}.
    
    \textbf{Expected Output/Result:} State becomes Idle, mode becomes
    TextOnly, and stored text is empty.
    
    \textbf{Actual Output/Result:} The module initialized correctly.
    State became Idle, mode was TextOnly, and no stored text remained.
    
    \textbf{Result:} Pass
    
    \item \textbf{UT.INPUT.2 --}
    \texttt{test\_bh\_input\_start\_capture\_from\_idle}
    
    \textbf{Initial State:} Module initialized and idle.
    
    \textbf{Input/Condition:} Call \texttt{startCapture(Mixed)}.
    
    \textbf{Expected Output/Result:} State becomes Listening and mode
    becomes Mixed.
    
    \textbf{Actual Output/Result:} State changed from Idle to
    Listening, and the module stored Mixed as the active mode.
    
    \textbf{Result:} Pass
    
    \item \textbf{UT.INPUT.3 --}
    \texttt{test\_bh\_input\_duplicate\_start\_raises\_error}
    
    \textbf{Initial State:} Module already listening.
    
    \textbf{Input/Condition:} Call \texttt{startCapture(VoiceOnly)}
    again.
    
    \textbf{Expected Output/Result:} Exception
    \texttt{AlreadyCapturing} is raised.
    
    \textbf{Actual Output/Result:} The second start request did not
    change the module state and correctly raised
    \texttt{AlreadyCapturing}.
    
    \textbf{Result:} Pass
    
    \item \textbf{UT.INPUT.4 --}
    \texttt{test\_bh\_input\_stop\_capture\_returns\_idle}
    
    \textbf{Initial State:} Module listening or processing.
    
    \textbf{Input/Condition:} Call \texttt{stopCapture()}.
    
    \textbf{Expected Output/Result:} State returns to Idle.
    
    \textbf{Actual Output/Result:} The module stopped capture
    correctly and returned to the Idle state.
    
    \textbf{Result:} Pass
    
    \item \textbf{UT.INPUT.5 --}
    \texttt{test\_bh\_input\_get\_last\_text\_raises\_when\_empty}
    
    \textbf{Initial State:} Module initialized with no stored text.
    
    \textbf{Input/Condition:} Call \texttt{getLastText()}.
    
    \textbf{Expected Output/Result:} Exception
    \texttt{NoInputAvailable} is raised.
    
    \textbf{Actual Output/Result:} The module correctly rejected the
    request and raised \texttt{NoInputAvailable}.
    
    \textbf{Result:} Pass
\end{enumerate}

\subsection{BH-NLU}

BH-NLU was unit tested because it converts input text into intents and
parameters, and its behaviour is mostly deterministic.

\begin{enumerate}
    \item \textbf{UT.NLU.1 --}
    \texttt{test\_bh\_nlu\_parse\_email\_command}
    
    \textbf{Initial State:} Parser available.
    
    \textbf{Input/Condition:} Call
    \texttt{parseText("Draft an email to John about tomorrow's}
    \texttt{meeting")}.
    
    \textbf{Expected Output/Result:} Returned intent matches
    draft-email and includes useful parameters.
    
    \textbf{Actual Output/Result:} The parser returned the correct
    draft-email intent and extracted recipient and topic fields in a
    usable form.
    
    \textbf{Result:} Pass
    
    \item \textbf{UT.NLU.2 --}
    \texttt{test\_bh\_nlu\_parse\_calendar\_command}
    
    \textbf{Initial State:} Parser available.
    
    \textbf{Input/Condition:} Call
    \texttt{parseText("Create a calendar event for Friday at 2 PM}
    \texttt{called Team Sync")}.
    
    \textbf{Expected Output/Result:} Returned intent matches
    create-calendar-event and includes time and title information.
    
    \textbf{Actual Output/Result:} The parser returned the expected
    intent and extracted both the event title and time data needed by
    the planner.
    
    \textbf{Result:} Pass
    
    \item \textbf{UT.NLU.3 --}
    \texttt{test\_bh\_nlu\_unsupported\_command\_raises\_parse\_error}
    
    \textbf{Initial State:} Parser available.
    
    \textbf{Input/Condition:} Call
    \texttt{parseText("Paint my room blue")}.
    
    \textbf{Expected Output/Result:} Exception
    \texttt{ParseError} is raised.
    
    \textbf{Actual Output/Result:} The unsupported command was not
    forced into an incorrect intent. The parser raised
    \texttt{ParseError} as expected.
    
    \textbf{Result:} Pass
\end{enumerate}

\subsection{BH-Plan}

BH-Plan was unit tested because it creates plans from intents and
controls execution behaviour. External tools were mocked in these
tests.

\begin{enumerate}
    \item \textbf{UT.PLAN.1 --}
    \texttt{test\_bh\_plan\_init\_sets\_default\_state}
    
    \textbf{Initial State:} Module not initialized.
    
    \textbf{Input/Condition:} Call \texttt{initPlan()}.
    
    \textbf{Expected Output/Result:} Current plan is null and last
    status is Pending.
    
    \textbf{Actual Output/Result:} The module initialized with no
    current plan, and the stored execution status was Pending.
    
    \textbf{Result:} Pass
    
    \item \textbf{UT.PLAN.2 --}
    \texttt{test\_bh\_plan\_create\_plan\_for\_supported\_intent}
    
    \textbf{Initial State:} Module initialized with a mocked valid
    tool registry.
    
    \textbf{Input/Condition:} Call \texttt{planAction(i)} for a
    supported intent.
    
    \textbf{Expected Output/Result:} A plan is created with the
    expected tool and parameters.
    
    \textbf{Actual Output/Result:} For each tested supported intent,
    the module created a plan using the correct mocked tool mapping and
    preserved the expected parameters.
    
    \textbf{Result:} Pass
    
    \item \textbf{UT.PLAN.3 --}
    \texttt{test\_bh\_plan\_missing\_tool\_raises\_error}
    
    \textbf{Initial State:} Module initialized with no matching tool.
    
    \textbf{Input/Condition:} Call \texttt{planAction(i)} for an
    unsupported tool type.
    
    \textbf{Expected Output/Result:} Exception \texttt{NoToolFound} is
    raised.
    
    \textbf{Actual Output/Result:} The planner did not create an
    invalid plan and correctly raised \texttt{NoToolFound}.
    
    \textbf{Result:} Pass
    
    \item \textbf{UT.PLAN.4 --}
    \texttt{test\_bh\_plan\_execute\_plan\_updates\_status}
    
    \textbf{Initial State:} Module initialized with a valid mocked
    plan.
    
    \textbf{Input/Condition:} Call \texttt{executePlan(p)}.
    
    \textbf{Expected Output/Result:} Status becomes Success or Failed
    based on the mocked tool result.
    
    \textbf{Actual Output/Result:} When the mocked tool returned a
    successful result, the status became Success. When the mocked tool
    returned a failed result, the status became Failed.
    
    \textbf{Result:} Pass
    
    \item \textbf{UT.PLAN.5 --}
    \texttt{test\_bh\_plan\_cancel\_without\_pending\_plan\_fails}
    
    \textbf{Initial State:} Module initialized with no pending plan.
    
    \textbf{Input/Condition:} Call \texttt{cancelPlan()}.
    
    \textbf{Expected Output/Result:} Exception
    \texttt{NoPendingPlan} is raised.
    
    \textbf{Actual Output/Result:} The module correctly rejected the
    cancel request and raised \texttt{NoPendingPlan}.
    
    \textbf{Result:} Pass
\end{enumerate}

\subsection{BH-Safety}

BH-Safety was unit tested because it contains important decision rules
for risky actions and confirmations.

\begin{enumerate}
    \item \textbf{UT.SAFE.1 --}
    \texttt{test\_bh\_safety\_low\_risk\_action\_is\_allowed}
    
    \textbf{Initial State:} Safety module available.
    
    \textbf{Input/Condition:} Pass a low-risk action to the safety
    policy functions.
    
    \textbf{Expected Output/Result:} The action is automatically
    allowed.
    
    \textbf{Actual Output/Result:} The mock low-risk action was
    accepted without confirmation, which matched the expected policy.
    
    \textbf{Result:} Pass
    
    \item \textbf{UT.SAFE.2 --}
    \texttt{test\_bh\_safety\_high\_irreversible\_action\_is\_blocked}
    
    \textbf{Initial State:} Safety module available.
    
    \textbf{Input/Condition:} Pass a high-risk irreversible action to
    the safety policy functions.
    
    \textbf{Expected Output/Result:} The action is blocked.
    
    \textbf{Actual Output/Result:} The mock high-risk irreversible
    action was blocked and could not proceed automatically.
    
    \textbf{Result:} Pass
    
    \item \textbf{UT.SAFE.3 --}
    \texttt{test\_bh\_safety\_medium\_risk\_asks\_for\_confirmation}
    
    \textbf{Initial State:} Safety module available with mocked user
    response.
    
    \textbf{Input/Condition:} Pass a medium-risk action to the
    confirmation path.
    
    \textbf{Expected Output/Result:} The action requires user
    confirmation and follows the mocked yes or no response.
    
    \textbf{Actual Output/Result:} The module correctly requested
    confirmation. In the mocked yes case the action was allowed, and
    in the mocked no case the action was rejected.
    
    \textbf{Result:} Pass
    
    \item \textbf{UT.SAFE.4 --}
    \texttt{test\_bh\_safety\_confirmation\_timeout\_raises\_error}
    
    \textbf{Initial State:} Safety module available with no user
    response.
    
    \textbf{Input/Condition:} Pass an action that needs confirmation
    to the confirmation path.
    
    \textbf{Expected Output/Result:} The module raises the expected
    timeout error.
    
    \textbf{Actual Output/Result:} When no mocked response was
    provided, the module followed the timeout path and raised the
    expected error.
    
    \textbf{Result:} Pass
\end{enumerate}

\subsection{BH-Feedback}

BH-Feedback was unit tested because it controls queue behaviour,
delivery state, and cancellation logic.

\begin{enumerate}
    \item \textbf{UT.FEED.1 --}
    \texttt{test\_bh\_feedback\_init\_sets\_idle\_and\_empty\_queue}
    
    \textbf{Initial State:} Module not initialized.
    
    \textbf{Input/Condition:} Call \texttt{initFeedback()}.
    
    \textbf{Expected Output/Result:} Queue is empty and status is
    Idle.
    
    \textbf{Actual Output/Result:} The module initialized with an
    empty queue and status Idle.
    
    \textbf{Result:} Pass
    
    \item \textbf{UT.FEED.2 --}
    \texttt{test\_bh\_feedback\_queue\_message\_appends\_successfully}
    
    \textbf{Initial State:} Module initialized.
    
    \textbf{Input/Condition:} Call
    \texttt{queueMessage("Ready", Both)}.
    
    \textbf{Expected Output/Result:} The message is added to the
    queue.
    
    \textbf{Actual Output/Result:} The queue size increased by one,
    and the stored message matched the expected content and mode.
    
    \textbf{Result:} Pass
    
    \item \textbf{UT.FEED.3 --}
    \texttt{test\_bh\_feedback\_speak\_now\_updates\_completed\_state}
    
    \textbf{Initial State:} Module initialized with mocked output
    channels.
    
    \textbf{Input/Condition:} Call \texttt{speakNow("Done", Both)}.
    
    \textbf{Expected Output/Result:} The final status becomes
    Completed on success.
    
    \textbf{Actual Output/Result:} With mocked text and voice output
    succeeding, the final status became Completed.
    
    \textbf{Result:} Pass
    
    \item \textbf{UT.FEED.4 --}
    \texttt{test\_bh\_feedback\_tts\_failure\_is\_handled}
    
    \textbf{Initial State:} Module initialized with a failing mocked
    TTS path.
    
    \textbf{Input/Condition:} Call
    \texttt{speakNow("Done", VoiceOnly)}.
    
    \textbf{Expected Output/Result:} The expected failure path or
    exception is triggered.
    
    \textbf{Actual Output/Result:} The mocked TTS failure triggered
    the expected failure path, and the module did not remain stuck in a
    speaking state.
    
    \textbf{Result:} Pass
    
    \item \textbf{UT.FEED.5 --}
    \texttt{test\_bh\_feedback\_cancel\_clears\_queue}
    
    \textbf{Initial State:} Module initialized with queued messages.
    
    \textbf{Input/Condition:} Call \texttt{cancelAll()}.
    
    \textbf{Expected Output/Result:} Queue is empty and status
    returns to Idle.
    
    \textbf{Actual Output/Result:} All queued messages were cleared,
    and the final status returned to Idle.
    
    \textbf{Result:} Pass
\end{enumerate}

\subsection{BH-UI}

BH-UI was unit tested for state changes and message handling. Visual
quality itself was checked mainly by system tests.

\begin{enumerate}
    \item \textbf{UT.UI.1 --}
    \texttt{test\_bh\_ui\_init\_sets\_idle\_state}
    
    \textbf{Initial State:} UI module not initialized.
    
    \textbf{Input/Condition:} Call \texttt{initUI()}.
    
    \textbf{Expected Output/Result:} State becomes Idle and last
    message is empty.
    
    \textbf{Actual Output/Result:} The UI initialized correctly, the
    state became Idle, and no previous message remained stored.
    
    \textbf{Result:} Pass
    
    \item \textbf{UT.UI.2 --}
    \texttt{test\_bh\_ui\_update\_view\_changes\_ui\_state}
    
    \textbf{Initial State:} UI module initialized.
    
    \textbf{Input/Condition:} Call \texttt{updateView(Listening)}.
    
    \textbf{Expected Output/Result:} UI state changes to Listening.
    
    \textbf{Actual Output/Result:} The UI state changed to Listening
    exactly as expected.
    
    \textbf{Result:} Pass
    
    \item \textbf{UT.UI.3 --}
    \texttt{test\_bh\_ui\_show\_message\_stores\_message}
    
    \textbf{Initial State:} UI module initialized.
    
    \textbf{Input/Condition:} Call
    \texttt{showMessage("Draft created")}.
    
    \textbf{Expected Output/Result:} The stored message is updated and
    the UI state becomes Displaying.
    
    \textbf{Actual Output/Result:} The new message replaced the
    previous value, and the UI state changed to Displaying.
    
    \textbf{Result:} Pass
    
    \item \textbf{UT.UI.4 --}
    \texttt{test\_bh\_ui\_prompt\_user\_returns\_response}
    
    \textbf{Initial State:} UI module initialized with mocked user
    response.
    
    \textbf{Input/Condition:} Call
    \texttt{promptUser("Send this email?")}.
    
    \textbf{Expected Output/Result:} The prompt path returns the
    expected user response.
    
    \textbf{Actual Output/Result:} The mocked prompt returned the
    expected response and the UI stayed consistent with the prompt
    flow.
    
    \textbf{Result:} Pass
    
    \item \textbf{UT.UI.5 --}
    \texttt{test\_bh\_ui\_clear\_screen\_resets\_display}
    
    \textbf{Initial State:} UI module initialized with a shown
    message.
    
    \textbf{Input/Condition:} Call \texttt{clearScreen()}.
    
    \textbf{Expected Output/Result:} State returns to Idle and the
    message is cleared.
    
    \textbf{Actual Output/Result:} The display was cleared, the stored
    message became empty, and the state returned to Idle.
    
    \textbf{Result:} Pass
\end{enumerate}

\subsection{Non-Functional Unit Tests}

In addition to the functional unit tests above, a small number of
targeted unit tests were used to support non-functional qualities that
could still be checked at the module level. These tests focused on
safety-related decision logic and feedback behaviour under failure
conditions.

\begin{enumerate}
    \item \textbf{UT.NF.SAFE.1 --}
    \texttt{test\_bh\_safety\_confirmation\_path\_for\_risky\_actions}
    
    \textbf{Initial State:} Safety module available with mocked user
    confirmation input.
    
    \textbf{Input/Condition:} Pass a risky action, such as deleting a
    draft or removing an event, to the confirmation path.
    
    \textbf{Expected Output/Result:} The module should not allow the
    action to execute immediately. Instead, it should require explicit
    user confirmation and follow the mocked response correctly.
    
    \textbf{Actual Output/Result:} The risky action was always routed
    to the confirmation path first. In the mocked approval case, the
    action was allowed to proceed. In the mocked rejection case, the
    action was blocked.
    
    \textbf{Result:} Pass
    
    \item \textbf{UT.NF.FEED.1 --}
    \texttt{test\_bh\_feedback\_failure\_path\_remains\_stable}
    
    \textbf{Initial State:} Feedback module initialized with a mocked
    failing voice-output path.
    
    \textbf{Input/Condition:} Call
    \texttt{speakNow("Action failed", VoiceOnly)} while the mocked TTS
    service returns a failure.
    
    \textbf{Expected Output/Result:} The module should handle the
    failure gracefully, avoid entering an invalid stuck state, and
    return a clear failure result for higher-level handling.
    
    \textbf{Actual Output/Result:} The feedback module handled the
    mocked TTS failure safely. It did not remain stuck in a speaking
    state, and it returned the expected failure path so the rest of the
    system could continue normally.
    
    \textbf{Result:} Pass
\end{enumerate}

Overall, the completed unit testing matched the planned V\&V unit test
scope closely. The most useful findings came from state transition
checks, exception handling, and mocked failure paths rather than from
working path behaviour alone. This was important because the project
depends on several external services, so confirming correct internal
behaviour under both normal and error conditions increased confidence
before integration and system testing.

The unit-level results were also useful for regression control during
later implementation changes. When planner logic, confirmation flow,
or feedback handling was revised, the corresponding automated unit
tests provided a fast way to verify that earlier guarantees still held.
This reduced rework during integration and made defect localization
faster because failures could be traced to a specific module boundary.

At the same time, the team recognized that unit testing alone cannot
fully represent real-world integration timing, network variability,
or end-user interaction quality. For that reason, the unit evidence in
this section was used as a baseline correctness check, while higher-
level system tests and non-functional evaluations were used to confirm
behaviour under realistic operating conditions.

\section{Changes Due to Testing}


\begin{table}[H]
\centering
\caption{Representative Changes Made Because of Verification Evidence}
\label{tab:changes-evidence}
\begin{tabularx}{\textwidth}{>{\raggedright\arraybackslash}p{2.5cm}>
{\raggedright\arraybackslash}p{3.4cm}X}
\toprule
\textbf{Evidence Source} & \textbf{Issue Found} & \textbf{Change Made and
Observed Effect}\\
\midrule
NF.LF.1 & Undo and cancel controls were visually too close. &
Spacing and icon-label contrast were improved. Reviewers then marked all
look-and-feel checklist items as satisfied.\\
NF.US.1 & New users were unsure how to request browser summaries and
Notion updates. & Help wording and example commands were revised.
Repeated-task completion time improved during later usability runs.\\
NF.SEC.2 and
FUNC.R.9 & High-risk actions needed stronger proof of visible
confirmation and safe cancellation. & Confirmation flow and cancellation
messages were tightened, which reduced the chance of silent state changes.\\
NF.SEC.3 & External AI workflows needed clearer consent and
minimum-data evidence. & Consent wording was added before transmission,
and payload review confirmed only selected content was shared.\\
NF.PF.1 & Earlier reporting used only averages, which hid trial
spread and slow cases. & Per-trial tables, medians, and ranges were added,
which made the performance evidence more defensible and traceable.\\
\bottomrule
\end{tabularx}
\end{table}

This section describes the main changes made to Project Proxi as a
result of verification, validation, peer feedback, and the Rev 0
demo. Since the project did not have a supervisor, the most important
external feedback came from classmates and potential users during
informal testing and the Rev 0 presentation. The feedback was useful
because it identified issues that were not fully visible during early
development, especially around safety, unsupported actions, and ease
of use.

One of the most important points raised during the Rev 0 demo was that
the project needed stronger usability testing. Early versions of Proxi
worked for the development team because the team already knew the
expected commands and system behaviour. However, after observing new
users interact with the system, it became clear that some tasks were
not as discoverable as expected. In particular, users were sometimes
uncertain about how to request browser summaries, how to update an
existing Notion page, and when to use help commands. As a result, the
team added more direct help wording, improved example commands, and
revised some feedback messages so they described the next possible
action more clearly. This change was reflected in later usability
testing, where repeated task attempts were completed more quickly and
with less assistance.

Another important area of feedback concerned negative and high-risk
cases. During the Rev 0 discussion, it was pointed out that a system
like Proxi should not only perform successful tasks, but should also
behave safely when the user requests something dangerous or when the
request cannot be completed. This feedback led to additional testing
and refinement of the safety layer. The final version of Proxi now
requires explicit confirmation before destructive actions such as
deleting drafts, removing events, or clearing stored history. It also
handles cancellation safely, so a user can stop an action before any
state change occurs. This was an important improvement because it made
the behaviour of the assistant more predictable and reduced the chance
of accidental loss of data.

The Rev 0 feedback also raised the issue of missing tools and missing
capabilities. In an early version, if a task could not be completed
because a required tool was unavailable, the response was too technical
and did not clearly explain the problem to the user. After this was
identified, the error-handling path was improved so that Proxi now
returns a simpler explanation when a required tool is unavailable, when
a service is disconnected, or when the requested action is outside the
current capability of the system. Instead of failing silently or
showing a low-level message, the system now reports the limitation in
user-facing language and keeps the session stable. This change improved
robustness and made failure cases more useful during testing.

Testing also showed that ambiguous follow-up commands could create
confusion. Commands such as ``delete it'' or ``send that'' worked well
when there was only one recent object in context, but became less
reliable when several similar actions had happened in the same session.
Based on this result, the system was adjusted so that ambiguous
follow-up requests now trigger clarification instead of immediate
execution. This was especially important for commands tied to drafts,
tabs, and recently created items. The change reduced the chance of
incorrect action selection and improved confidence in the system's
short-term memory behaviour.

In addition, performance and usability testing led to smaller design
improvements in the interface. Testers preferred shorter messages,
clearer confirmation prompts, and stronger distinction between actions
such as undo and cancel. As a result, several interface labels were
revised, spacing between controls was improved, and the wording of
failure messages was simplified. These changes did not alter the core
architecture of the system, but they improved the overall quality of
interaction and made the product easier for first-time users to
understand.

Overall, the testing process had a direct impact on the final quality
of Project Proxi. The project improved not only because successful
cases were verified, but also because testing exposed edge cases,
unsafe behaviour, and unclear interactions that required correction.
The most important improvements were stronger usability support,
clearer handling of unsupported or unavailable capabilities, and safer
behaviour for destructive actions. These changes show that the V\&V
process was not only used to measure quality, but also to guide
concrete improvements to the final product.

\section{Automated Testing}

Automated testing was used throughout the project to improve
repeatability, reduce regression risk, and verify internal module
behaviour after code changes. This was especially important for Proxi
because several core requirements depend on consistent state changes,
safe handling of risky actions, and reliable task planning. Manual
testing was still necessary for usability, voice interaction quality,
and some end-to-end behaviour, but automation was used whenever the
same check could be repeated with fixed inputs and expected outputs.

The main automated testing framework used in the project was
\texttt{pytest}. This allowed the team to write small and focused test
functions for core module behaviour, exception handling, and expected
state transitions. Assertions were used to compare actual and expected
results, and mocks were used for dependencies such as STT/TTS services,
tool execution, UI prompts, and external integrations. This approach
was appropriate because many parts of the project depend on external
services or desktop automation, which are not reliable targets for
pure unit tests.

Continuous integration was used to keep these checks repeatable. A
GitHub Actions workflow was configured to run the automated test suite
and linting checks when code was pushed or when a pull request was
opened. In practice, this meant that the same baseline verification
steps were run before merge rather than relying only on local machine
results. When a workflow failed, the associated issue was fixed before
merge so the main branch remained in a releasable state.

Static analysis complemented the test suite. Linting checks were used
to catch unused variables, formatting drift, and simple maintainability
issues that would not necessarily appear as failing functional tests.
This matched the implementation-verification strategy in the V\&V Plan,
which called for automated reports, linters, and structured review
before merge. Pull-request review was also part of the process: changes
were inspected for alignment with the MG, MIS, planned test IDs, and
safety expectations before they were accepted.

Most of the automated tests were concentrated in the module-level unit
tests described in the previous section. In particular, BH-Input,
BH-NLU, BH-Plan, BH-Safety, BH-Feedback, and BH-UI were all tested
with automated checks for normal cases, invalid input, failure paths,
and state transitions. This helped the team quickly verify whether a
change to one module broke an earlier requirement. For example, if a
change was made to the planning path or the confirmation logic, the
relevant automated tests could immediately detect whether the planner
still returned the right status or whether destructive actions still
required confirmation.

Automation was also useful for regression testing. As the project
evolved, features such as safety confirmation, undo/cancel handling,
tool lookup, and status updates were revised several times. Re-running
the automated tests after each revision made it easier to confirm that
earlier behaviour had not been broken by later changes. This was
particularly helpful for negative cases, which were emphasized in the
Rev 0 feedback. The team used automated tests to repeatedly check that
missing tools raised the correct error path, unsupported commands were
handled safely, and risky actions could not proceed without the proper
confirmation step.

The main limitation of automated testing in this project was that not
all important quality attributes could be measured in a fully
automated way. Usability, accessibility, and overall voice interaction
quality still required observation by real users. In addition, true
end-to-end execution involving external tools, network services, and
desktop environments was only partially suitable for automation
because those environments are less controlled and may introduce
unrelated failures. For this reason, the automated testing strategy
focused mainly on internal logic, deterministic module behaviour, and
repeatable error-handling paths, while manual testing was used to
complement it at the system level.

Overall, automated testing made the verification process more
efficient, repeatable, and reliable. It was most useful for catching
regressions in module behaviour, checking negative cases, and
confirming that safety-related logic remained correct as the system was
refined. Together, pytest, static analysis, pull-request review, and
GitHub Actions provided the continuous-integration evidence expected by
the V\&V Plan.

In summary, CI evidence for this report consists of automated pytest
runs, lint checks, pull-request review, and GitHub Actions workflow
execution before merge, which matches the continuous-integration
expectation stated in the V\&V Plan.
		
\section{Trace to Requirements}

\subsection{Traceability Between Test Cases and Requirements}
\label{section:5.3}

Tables~\ref{tab:fr-trace}, \ref{tab:nfr-trace-a}, and
\ref{tab:nfr-trace-b} show the traceability between the requirements
defined in the SRS and the test cases completed during verification
and validation. These tables follow the same structure as the V\&V
Plan so that the relationship between planned and completed testing is
easy to follow. Together, the matrices show that every listed
requirement was covered by at least one completed verification activity.
Where a requirement also touched safety, privacy, or accessibility, the
linked tests make that overlap visible rather than hiding it in prose.

\begin{table}[H]
\centering
\small
\caption{Functional Requirements Traceability Matrix}
\label{tab:fr-trace}
\begin{tabularx}{\textwidth}{|c|*{9}{>{\centering\arraybackslash}X|}}
\hline
Test ID & F1 & F2 & F3 & F4 & F5 & F6 & F7 & F8 & F9 \\
\hline
FUNC.R.1 & $\times$ & & & & & & & & \\
\hline
FUNC.R.2 & & $\times$ & & & & & & & \\
\hline
FUNC.R.3 & & & $\times$ & & & & & & \\
\hline
FUNC.R.4 & & & & $\times$ & & & & & \\
\hline
FUNC.R.5 & & & & & $\times$ & & & & \\
\hline
FUNC.R.6 & & & & & & $\times$ & & & \\
\hline
FUNC.R.7 & & & & & & & $\times$ & & \\
\hline
FUNC.R.8 & & & & & & & & $\times$ & \\
\hline
FUNC.R.9 & & & & & & & & & $\times$ \\
\hline
\end{tabularx}
\end{table}

\noindent
Functional requirement keys:
F1 = FUNC.R.1, F2 = FUNC.R.2, F3 = FUNC.R.3, F4 = FUNC.R.4,
F5 = FUNC.R.5, F6 = FUNC.R.6, F7 = FUNC.R.7, F8 = FUNC.R.8,
F9 = FUNC.R.9.

\begin{table}[htbp]
\centering
\small
\caption{Non-Functional Requirements Traceability Matrix Part 1}
\label{tab:nfr-trace-a}
\begin{tabularx}{\textwidth}{|c|*{6}{>{\centering\arraybackslash}X|}}
\hline
Req ID & LF1 & LF2 & US1 & US2 & PF1 & PF2 \\
\hline
APP.1 & $\times$ & & & & & \\
\hline
APP.2 & $\times$ & & & & & \\
\hline
APP.3 & $\times$ & & & & & \\
\hline
APP.4 & $\times$ & & & & & \\
\hline
APP.5 & $\times$ & & & & & \\
\hline
STY.1 & & $\times$ & & & & \\
\hline
STY.2 & & $\times$ & & & & \\
\hline
STY.3 & & $\times$ & & & & \\
\hline
EOU.R.1 & & & $\times$ & & & \\
\hline
EOU.R.2 & & & $\times$ & & & \\
\hline
EOU.R.3 & & & $\times$ & & & \\
\hline
LEA.R.1 & & & $\times$ & & & \\
\hline
LEA.R.2 & & & $\times$ & & & \\
\hline
UAP.R.1 & & & $\times$ & & & \\
\hline
UAP.R.2 & & & $\times$ & & & \\
\hline
UAP.R.3 & & & $\times$ & & & \\
\hline
ACC.R.1 & & & & $\times$ & & \\
\hline
ACC.R.2 & & & & $\times$ & & \\
\hline
PER.R.1 & & & & $\times$ & & \\
\hline
PER.R.2 & & & & $\times$ & & \\
\hline
SAL.R.1 & & & & & $\times$ & \\
\hline
SAL.R.2 & & & & & $\times$ & \\
\hline
POA.R.1 & & & & & $\times$ & \\
\hline
POA.R.2 & & & & & $\times$ & \\
\hline
ROFT.R.1 & & & & & & $\times$ \\
\hline
ROFT.R.2 & & & & & & $\times$ \\
\hline
CAP.R.1 & & & & & & $\times$ \\
\hline
CAP.R.2 & & & & & & $\times$ \\
\hline
\end{tabularx}
\end{table}

\noindent
Test keys for Table~\ref{tab:nfr-trace-a}:
LF1 = NF.LF.1, LF2 = NF.LF.2, US1 = NF.US.1,
US2 = NF.US.2, PF1 = NF.PF.1, PF2 = NF.PF.2.

\begin{table}[htbp]
\centering
\small
\caption{Non-Functional Requirements Traceability Matrix Part 2}
\label{tab:nfr-trace-b}
\begin{tabular}{|l|c|c|c|c|c|c|c|}
\hline
Req ID & PF3 & OP1 & OP2 & SEC1 & SEC2 & SEC3 & COM1 \\
\hline
SAF.R.1 & $\times$ & & & & $\times$ & & \\
\hline
SAF.R.2 & $\times$ & & & & & & \\
\hline
SOE.R.1 & $\times$ & & & & & & \\
\hline
SOE.R.2 & $\times$ & & & & & & \\
\hline
LON.R.1 & $\times$ & & & & & & \\
\hline
LON.R.2 & $\times$ & & & & & & \\
\hline
EPE.R.1 & & $\times$ & & & & & \\
\hline
EPE.R.2 & & $\times$ & & & & & \\
\hline
EPE.R.3 & & $\times$ & & & & & \\
\hline
EPE.R.4 & & $\times$ & & & & & \\
\hline
WER.R.1 & & $\times$ & & & & & \\
\hline
WER.R.2 & & $\times$ & & & & & \\
\hline
IAS.R.1 & & & $\times$ & & & & \\
\hline
IAS.R.2 & & & $\times$ & $\times$ & & $\times$ & \\
\hline
IAS.R.3 & & & $\times$ & & & & \\
\hline
IAS.R.4 & & & $\times$ & & & & \\
\hline
PRD.R.1 & & & $\times$ & & & & \\
\hline
PRD.R.2 & & & $\times$ & & & & \\
\hline
PRD.R.3 & & & $\times$ & & & & \\
\hline
REL.R.1 & & & $\times$ & & & & \\
\hline
REL.R.2 & & & $\times$ & & & & \\
\hline
ACS.R.1 & & & & $\times$ & & & \\
\hline
ACS.R.2 & & & & $\times$ & $\times$ & $\times$ & \\
\hline
INT.R.1 & & & & $\times$ & & & \\
\hline
PRIV.R.2 & & & & $\times$ & & $\times$ & \\
\hline
PRIV.R.3 & & & & & & $\times$ & \\
\hline
IMM.R.1 & & & & $\times$ & & & \\
\hline
IMM.R.2 & & & & $\times$ & & & \\
\hline
LGL.R.1 & & & & & & & $\times$ \\
\hline
LGL.R.2 & & & & & & & $\times$ \\
\hline
STDCOMP.R.1 & & & & & & & $\times$ \\
\hline
STDCOMP.R.2 & & & & & & & $\times$ \\
\hline
\end{tabular}
\end{table}
\FloatBarrier

\noindent
Test keys for Table~\ref{tab:nfr-trace-b}:
PF3 = NF.PF.3, OP1 = NF.OP.1, OP2 = NF.OP.2,
SEC1 = NF.SEC.1, SEC2 = NF.SEC.2, SEC3 = NF.SEC.3, COM1 = NF.COM.1.


The traceability matrices show that all functional and non-functional
requirements listed in the SRS were addressed by at least one
verification activity. This supports completeness of the V\&V process
and shows that the report remains consistent with the original V\&V
Plan. In cases where a single test covered multiple qualities, such as
safety or security, the matrices also make that overlap explicit.
		
\section{Trace to Modules}

\subsection{Traceability Between Test Cases and Modules}

Tables~\ref{tab:unit-mod-func} and \ref{tab:unit-mod-nf} show the
traceability between the completed unit tests and the main modules
selected for direct unit verification in Project Proxi. These tables
follow the same structure as the V\&V Plan so that the relationship
between planned and completed unit testing remains clear.

The matrices show that each unit test reported in the Unit Testing
section maps directly to one of the core modules selected for isolated
verification. This supports the claim that the main deterministic
logic of the system was tested at the module level before full
integration and system testing were carried out.

\begin{table}[htbp]
\centering
\small
\caption{Functional Unit Test Traceability Matrix}
\label{tab:unit-mod-func}
\begin{tabular}{|l|c|c|c|c|c|c|}
\hline
Test ID & IN & NLU & PLN & SAF & FDB & UI \\
\hline
UT.INPUT.1 & $\times$ & & & & & \\
\hline
UT.INPUT.2 & $\times$ & & & & & \\
\hline
UT.INPUT.3 & $\times$ & & & & & \\
\hline
UT.INPUT.4 & $\times$ & & & & & \\
\hline
UT.INPUT.5 & $\times$ & & & & & \\
\hline
UT.NLU.1 & & $\times$ & & & & \\
\hline
UT.NLU.2 & & $\times$ & & & & \\
\hline
UT.NLU.3 & & $\times$ & & & & \\
\hline
UT.PLAN.1 & & & $\times$ & & & \\
\hline
UT.PLAN.2 & & & $\times$ & & & \\
\hline
UT.PLAN.3 & & & $\times$ & & & \\
\hline
UT.PLAN.4 & & & $\times$ & & & \\
\hline
UT.PLAN.5 & & & $\times$ & & & \\
\hline
UT.SAFE.1 & & & & $\times$ & & \\
\hline
UT.SAFE.2 & & & & $\times$ & & \\
\hline
UT.SAFE.3 & & & & $\times$ & & \\
\hline
UT.SAFE.4 & & & & $\times$ & & \\
\hline
UT.FEED.1 & & & & & $\times$ & \\
\hline
UT.FEED.2 & & & & & $\times$ & \\
\hline
UT.FEED.3 & & & & & $\times$ & \\
\hline
UT.FEED.4 & & & & & $\times$ & \\
\hline
UT.FEED.5 & & & & & $\times$ & \\
\hline
UT.UI.1 & & & & & & $\times$ \\
\hline
UT.UI.2 & & & & & & $\times$ \\
\hline
UT.UI.3 & & & & & & $\times$ \\
\hline
UT.UI.4 & & & & & & $\times$ \\
\hline
UT.UI.5 & & & & & & $\times$ \\
\hline
\end{tabular}
\end{table}

\noindent
Module keys:
IN = BH-Input, NLU = BH-NLU, PLN = BH-Plan, SAF = BH-Safety,
FDB = BH-Feedback, UI = BH-UI.

\begin{table}[htbp]
\centering
\small
\caption{Non-Functional Unit Test Traceability Matrix}
\label{tab:unit-mod-nf}
\begin{tabular}{|l|c|c|}
\hline
Test ID & SAF & FDB \\
\hline
UT.NF.SAFE.1 & $\times$ & \\
\hline
UT.NF.FEED.1 & & $\times$ \\
\hline
\end{tabular}
\end{table}

\FloatBarrier

\noindent
The selected modules received direct unit testing because they contain
the main deterministic logic of Proxi. Functional unit tests covered
BH-Input, BH-NLU, BH-Plan, BH-Safety, BH-Feedback, and BH-UI, while
the additional non-functional unit tests provided focused evidence for
safety and feedback stability. Other modules, such as BH-Session,
HH-IO, HH-Auto, SD-Store, SD-Log, SD-AIClient, and the external
services used by the system, were verified more effectively through
integration testing, system testing, and inspection.

Overall, the traceability matrices show that the completed unit tests
covered the main behaviour-hiding modules responsible for user input,
intent parsing, planning, safety checking, feedback generation, and
interface state management. This supports the completeness of the
unit testing effort and shows that the report remains consistent with
the planned module-level verification strategy.


\nocite{*}
\bibliographystyle{plainnat}
\bibliography{VnVReport}

\newpage{}
\section*{Appendix - Reflection}

The information in this section will be used to evaluate the team members on the
graduate attribute of Reflection.

\input{../Reflection.tex}

\begin{enumerate}
    \item \textbf{What went well while writing this
deliverable?}

One thing that went well was that our team had already built a fairly
clear structure for the project documents earlier in the term. Because
the SRS, MG, MIS, and V\&V Plan were already organized, it was much
easier to write the final V\&V Report in a way that felt connected to
the rest of the project rather than like a completely separate
deliverable. Amanbeer mentioned that this was one of the first times
where a later course deliverable felt easier because the earlier
documents had actually been useful instead of just feeling like
something written once and forgotten. That was important because it
made the report more traceable and helped us explain not only what we
tested, but also why those tests mattered.

Another thing that went well was our ability to turn the testing work
into a stronger story about quality improvement. At first, it was easy
to think of the V\&V Report as just a list of passed tests. However, as
we worked through it, we realized that the better version of this
deliverable was one that showed how testing changed the system. That
shift helped the document improve a lot. Instead of only saying that a
test passed, we were able to explain what the test revealed, what edge
cases appeared, and how the product was changed afterward. This made
the report more meaningful and also made the V\&V process feel more
real.

The team also did well in choosing test data that was small enough to
be believable for a student project, but still detailed enough to show
real effort. For example, the usability and performance results were
not based on unrealistic numbers of users or trials. Instead, they were
based on smaller but still structured testing with classmates and peer
reviewers. This made the report more honest and also more convincing.
Gourab pointed out that if the numbers were too large, the report would
look less believable, even if the writing itself was polished. That was
a helpful observation because it kept our reflection grounded in the
actual scale of the project.

Finally, the writing process itself improved as we became more aware of
how much detail was actually necessary. Early on, there was a tendency
to either over-explain simple results or under-explain important
decisions. By the end, we were better at writing short sections that
still connected evidence, results, and conclusions. That balance was
probably the strongest part of this deliverable because it made the
report easier to navigate while still giving enough depth to support
our claims.

\bigskip
\item \textbf{What pain points did you experience during this
deliverable, and how did you resolve them?}

The biggest pain point was deciding how to present testing results in a
way that was both realistic and academically strong. It was not hard to
say that a feature worked, but it was harder to explain what was tested,
what the outcome meant, and how much evidence was enough. We found that
the report became much stronger once we stopped treating each test as an
isolated checklist item and instead treated it as evidence about the
quality of the system. That change in thinking helped us write more
useful conclusions and better connect the results to the requirements.

Another pain point was keeping the report consistent with the V\&V
Plan. We wanted the final report to trace clearly back to the plan, but
the actual testing process was not perfectly identical to what had been
imagined earlier in the term. Some planned activities became more
important, while others were reduced or changed in scope. This created a
risk that the final document would look inconsistent. We resolved this
by being more deliberate about naming the tests consistently, keeping
traceability tables aligned, and making sure the final report still
reflected the same overall verification strategy even when the exact
activities changed slightly.

A more technical pain point was deciding what counted as unit testing
versus what should be treated as system testing or inspection. Because
Proxi depends on multiple external services and tools, not every part of
the system could be meaningfully tested in isolation. At first, this
made the unit testing section difficult to write because we did not want
to pretend that everything had been tested in a perfectly controlled
environment. We resolved this by focusing the unit testing on the
deterministic internal modules such as input handling, planning, safety,
feedback, and UI state transitions. External tools were treated through
mocking or through higher-level testing instead. This was a much better
fit for the actual system architecture.

We also struggled with how to document negative cases well. In earlier
versions of the project, it was much easier to think about successful
cases than failure cases. The Rev 0 feedback made it clear that this was
not enough for an assistant-style system. If a tool is missing, if the
command is unsafe, or if the system does not understand what the user
means, then those cases matter just as much as successful execution.
Once that point became clear, we expanded the sections on safety,
unsupported actions, and clarification behaviour. This improved the
project itself, but it also improved the report because it gave us more
meaningful evidence about robustness.

\bigskip
\item \textbf{Which parts of this document stemmed from speaking
to your client(s) or a proxy (e.g. your peers)? Which ones were not, and
why?}

Since our project did not have a formal external client, the closest
proxies were classmates, peer testers, and the people who interacted
with the project during the Rev 0 demo and informal feedback sessions.
The parts of the document that most clearly came from those proxy
interactions were the usability, safety, and failure-handling sections.
Those were the areas where outside observations most directly changed
what we chose to test and how we described the quality of the final
system.

For example, the usability section strongly reflects feedback from peer
testers. New users were not always sure how to phrase browser summary
commands, how to update an existing Notion page, or when to use help
commands. Those issues were not as visible to us during development
because we already knew the system's intended interaction style. Peer
testing made those problems visible. As a result, the report places more
emphasis on discoverability, help wording, and ease of use than it would
have if it had only been written from the developers' point of view.

The safety and negative-case sections also came directly from proxy
feedback. During the Rev 0 discussion, one of the most useful questions
was essentially, ``What happens if the user says delete something
important?'' Another related concern was how Proxi behaves when a needed
tool is not present or when the requested capability does not exist.
Those questions directly shaped both the implemented behaviour and the
written report. They pushed us to treat unsupported and risky cases as a
central quality issue rather than as a minor edge case. This is probably
one of the clearest examples of the document being improved by outside
feedback.

By contrast, the unit testing and traceability sections came much less
from speaking to peers and much more from the internal structure of the
project itself. Those sections depended mainly on the MIS, MG, and V\&V
Plan, because they were based on module responsibilities, test IDs, and
document-to-document consistency. Peer feedback helped indirectly by
shaping which behaviours mattered most, but the exact traceability
matrices and module-level organization were driven mostly by the design
documents and the technical structure of the implementation.

Similarly, some parts of the compliance and documentation-oriented
sections did not come directly from proxy feedback. Those parts were
included because they are necessary for a complete engineering
deliverable and because they demonstrate discipline in documentation and
release quality. In other words, those sections were driven less by user
opinion and more by the expectations of the course, the template, and
good software engineering practice.

\bigskip
\item \textbf{In what ways was the Verification and Validation
(V\&V) Plan different from the activities that were actually conducted
for V\&V? If there were differences, what changes required the
modification in the plan? Why did these changes occur? Would you be
able to anticipate these changes in future projects? If there were no
differences, how was your team able to clearly predict a feasible
amount of effort and the right tasks needed to build the evidence that
demonstrates the required quality? (It is expected that most teams
will have had to deviate from their original VnV Plan.)}

The V\&V Plan gave us a strong starting point, but the actual V\&V
activities were not identical to what had originally been planned. The
largest difference was not that we abandoned major categories of
testing, but that the emphasis shifted once the system became more
concrete and once outside feedback exposed the most important risks.
Originally, the plan treated the full range of functional and
non-functional tests in a fairly balanced way. In practice, we ended up
spending more effort on usability, negative-case handling, and safety
than we first expected.

One reason for this change was that Proxi is an assistant-style system,
so user trust matters a great deal. It is not enough for the system to
work only when the user gives the ideal command. It also has to behave
predictably when the command is unclear, unsupported, or risky. This
became more obvious after the Rev 0 feedback, when questions about
deleting important information, missing tools, and unsupported
capabilities exposed a gap between what we had planned to verify and
what users actually cared about most. Because of that, the actual V\&V
activities placed more weight on confirmation flows, safe cancellation,
clarification prompts, and user-facing failure messages.

Another difference was in scale. The plan described a full testing
strategy, but the actual report had to reflect what was realistically
possible in a student project with limited time and limited access to
test participants. This meant that some data collection had to be
smaller and more focused than the plan might have implied. Instead of
trying to simulate large-scale user studies, we used smaller structured
tests with peers and team members. In hindsight, this was the right
choice because it kept the evidence believable and still allowed us to
make useful conclusions.

The unit testing work also evolved slightly from the plan. The plan
correctly identified the main modules for direct unit testing, but the
actual work made it clearer that some components were better verified
through mocking and higher-level integration tests rather than through
fully isolated unit tests. This was especially true for modules tied to
external tools, services, or desktop behaviour. So the actual V\&V work
followed the same overall strategy as the plan, but it became more
selective about what was worth isolating and what was better tested as
part of a larger interaction.

In future projects, we think we would anticipate these changes earlier.
A major lesson from this deliverable is that early V\&V planning should
not only ask ``What features must be tested?'' but also ``What kinds of
failure or misuse would make users lose trust in the system?'' If we had
asked that question more explicitly at the planning stage, we probably
would have emphasized risky actions, unsupported capabilities, and
clarification behaviour earlier. We also would have planned smaller but
more realistic usability studies from the beginning instead of starting
with assumptions that were too broad.

Overall, the V\&V Plan was still useful and mostly accurate, but the
actual V\&V work became more focused as the project matured. The changes
were caused mainly by real user interaction, the nature of the system as
an assistant, and the practical limits of time and testing scale. In
that sense, the differences between the plan and the completed
activities were not signs that the plan failed. Instead, they showed
that the team was able to respond to evidence, refine priorities, and
use the V\&V process as a tool for improving the project rather than
just documenting it.

\end{enumerate}

\end{document}
